\documentclass[10pt,a4paper]{amsart}
\usepackage[margin=19mm]{geometry}
\usepackage{amsmath,amssymb,mathtools}
\usepackage[scr=rsfs,cal=euler]{mathalfa}
\usepackage{xfrac}
\usepackage[unicode,hidelinks,bookmarks=false]{hyperref}
\newcommand{\ie}{i.e.\ }
\newcommand{\cf}{cf.\ }
\newcommand{\resp}{respectively\ }
\newcommand{\Subsection}{\subsection}
\providecommand{\nobreakdash}{\nobreak\hbox{-}}
\setlength{\emergencystretch}{2em}
\multlinegap0pt
\usepackage{bm}
\usepackage{xcolor}
\newtheorem{theorem}{Theorem}
\newtheorem{proposition}[theorem]{Proposition}
\newtheorem{definition}[theorem]{Definition}
\newcommand{\eps}{\varepsilon}
\renewcommand{\geq}{\geqslant}
\title[The limiting constant for the KMP-to-KPZ convergence (Proposition 6.7)]{Convergence of the KMP model to the KPZ equation: the limiting constant of the scaling (Proposition 6.7)}
\author{Guillaume Barraquand, Francesco Casini}
\date{Source: arXiv:2507.19222v3 (CC BY 4.0)}
\begin{document}
\maketitle

\noindent\emph{Source and licence.} Convergence of the KMP model to the KPZ equation. Version arXiv:2507.19222v3 (CC BY 4.0), distributed under the Creative Commons Attribution 4.0 International licence, as recorded in the arXiv licence metadata for this exact version and shown on the abstract page https://arxiv.org/abs/2507.19222v3. Full rights notice: licenses/arxiv-2507.19222-CC-BY-4.0.txt.\par\smallskip

\noindent\textit{Editorial scope.} Counted unit: the proposition printed as Proposition 6.7 of the article (source0.tex lines 1107--1112, source label \texttt{proposition-limitingConstant}), delivered together with the article's own complete original proof of it (source0.tex lines 1113--1270, one \texttt{proof} environment of 158 lines that computes the limit of the numerator and of the denominator separately and concludes the stated value with no gap). No mathematics is written, repaired or re-derived: statement and proof are the article's own text line for line, in the article's own notation ($\gamma_\eps$, $N^{\eps}$, $D^{\eps}$, $\pi^{\text{inv}}_{\alpha}$, $Z^{\eps}(n)$, $\mathbf{p}^{\eps}$, $\alpha$). Required context retained uncounted: the display (normalized-piInv) fixing the re-normalized invariant measure of the random walk in random environment (lines 1006--1011); the article's Definition 6.2, which introduces the discrete-time Markov process $Z^{\eps}$ whose covariance appears in the numerator (lines 1024--1026, source label \texttt{remark-Zn}); and the display (epsilon-constant) defining the constant $\gamma_\eps$ whose limit is the object of the counted proposition (lines 1044--1046). Every definition, numbered display and label of those lines is retained unchanged. Omitted from this package and not redistributed: the article's preamble (lines 1--138); its title page and abstract (lines 139--166); its introduction with the announcement of the KMP-to-KPZ convergence theorem (lines 167--219); its second section on the KMP model as a stochastic flow, its third section stating the main result, its fourth section verifying the hypotheses, its fifth section on convergence in the Skorokhod space and its sixth section on the computation of the noise variance, apart from the three retained displays (lines 220--1005, 1012--1023, 1027--1043, 1047--1106); the rest of its sixth section after the counted proof (lines 1271--1693); its bibliography (line 1694, generated by the biblatex auxiliary file shipped in the e-print); and its closing end-document line (line 1695). Nothing inside a retained excerpt is omitted or altered: every retained body is delivered exactly as the article's lines are written, so no omission receipt of any kind accompanies this package. Citations: the retained excerpts carry no citation command at all, so no bibliography entry of the article is transcribed and no reference is redistributed. Reference adaptation: none was needed and none was made. Every reference inside the delivered unit resolves inside it: the label epsilon-constant, eqref'ed by the counted statement; the label normalized-piInv, eqref'ed by the counted proof; the labels N-epsilon, telescopic, cov-z1, z1-denominator, numerator-approx-final and denominator-approx-final of the displays inside the counted proof; and the label of the counted statement itself together with the inner display (statement-proposition-gamma). No unresolved reference is produced. Printed slips kept literal: no printed word, symbol, hypothesis, number or colour command of the counted statement or of its proof was altered; in particular the retained proof keeps the article's own \texttt{\textbackslash color\{black\}} annotations inside its displays. Layout and class: the article's own class is \texttt{amsart} with the options a4paper, 11pt, oneside and reqno, and it loads among others mathtools, amssymb, bbm, bm, dsfont, tikz and hyperref; no publisher class or style file is used or shipped, so none travels with this package and none is redistributed. The wrapper is the lane's pinned \texttt{amsart} editorial wrapper at 10pt on A4, which adds the article's own macro definitions that the retained text needs ($\eps$ and the article's own slanted $\geq$) and loads the \texttt{bm} and \texttt{xcolor} packages that the retained proof needs; it declares the theorem-like environments the retained text needs (\texttt{theorem}, \texttt{proposition}, \texttt{definition}) with the article's own shared counter but sequentially, because the wrapper has no sections. The delivered numbering is the unit's own: the retained Definition 6.2 prints as Definition 1 and the counted Proposition 6.7 prints as Proposition 2. The article's own \texttt{\textbackslash numberwithin\{equation\}\{section\}} is not reproduced, so equations are numbered sequentially in the wrapper instead of by section; that and the 11pt-to-10pt change are page-appearance differences of the wrapper only, and no formula, symbol, hypothesis or argument changes. Assets: no retained span contains a figure environment, a graphics-inclusion command, a PDF-page inclusion, a table or a TikZ picture; the article's own figures and drawings live in the omitted part of the paper, no image file is copied into this package, the package declares no asset dependency and no image file is redistributed with it. Licence: version arXiv:2507.19222v3 of this article is distributed under the Creative Commons Attribution 4.0 International licence, as recorded in the arXiv licence metadata for this exact version and shown on the abstract page https://arxiv.org/abs/2507.19222v3. A full rights notice is delivered at licenses/arxiv-2507.19222-CC-BY-4.0.txt. Characters: every retained excerpt is pure ASCII (the article's twelve non-ASCII characters all lie outside the retained spans), so the delivered engine, XeLaTeX with its default Latin Modern text font, reports no missing character. No glyph is substituted and no word is rewritten.
\begin{equation}\label{normalized-piInv}
	\pi^{\text{inv}}_{\alpha}(x)=\begin{cases}
		1&\text{if}\quad x\neq 0\\
		\frac{(\alpha+1)}{\alpha}&\text{if}\quad x=0
	\end{cases}\,.
\end{equation}

\begin{definition}\label{remark-Zn}
	We denote by $(Z^{\epsilon}(n))_{n\in \mathbb{Z}_{\geq 0}}$ the discrete-time Markov process with transition kernel $\mathbf{p}_{\text{dif},n}^{\epsilon}(\cdot,\cdot)$. 
\end{definition}

\begin{equation}\label{epsilon-constant}
	\gamma_{\epsilon}^{2}:=\frac{N^{\epsilon}}{D^{\epsilon}}\,.
\end{equation}

\begin{proposition}\label{proposition-limitingConstant}
	Recall $\gamma_{\epsilon}$ from \eqref{epsilon-constant}. We have
	\begin{equation}\label{statement-proposition-gamma}
		\lim_{\epsilon\to 0}\gamma_{\epsilon}^{2}=\frac{1}{{\color{black}2}\alpha}\,.
	\end{equation}
\end{proposition}

\begin{proof} We treat the discrete-time numerator and denominator separately. 
	Let us start with the numerator. We re-write $N^{\epsilon}$ as
	\begin{equation}\label{N-epsilon}
		N^{\epsilon}=	\sum_{z\in \mathbb{Z}}\pi^{\text{inv}}_{\alpha}(z)\mathbb{COV}^{(\epsilon)}\left(\mathrm{E}_{z}^{(\epsilon)}\left[X^{\epsilon}(\lfloor \epsilon^{-1}\rfloor)-z\right],\mathrm{E}_{0}^{(\epsilon)}[Y^{\epsilon}(\lfloor \epsilon^{-1}\rfloor)]\right)\,.
	\end{equation}
		Above, we have denoted by $\mathrm{E}^{(\epsilon)}_{z}$ the quenched expectation with respect to the path space measure of the RWRE initialized at $z$. 
			Using the telescopic identity 
		\begin{equation}\label{telescopic}
			X^{\epsilon}(\lfloor\epsilon^{-1}\rfloor)-X^{\epsilon}(0)=\sum_{i=0}^{\lfloor\epsilon^{-1}\rfloor-1}\left(X^{\epsilon}(i+1)-X^{\epsilon}(i)\right)\,,
		\end{equation}
		and using the fact that the two walkers can be correlated only through the same Bernoulli event, we have that 
		\begin{equation*}
			\begin{split}
				N^{\epsilon}=&
				\sum_{i=0}^{\lfloor\epsilon^{-1}\rfloor-1}\sum_{z\in \mathbb{Z}}\pi^{\text{inv}}_{\alpha}(z)\mathbb{COV}^{(\epsilon)}\left(\mathrm{E}_{z}^{(\epsilon)}\left[X^{\epsilon}(i+1)-X^{\epsilon}( i)\right],\mathrm{E}_{0}^{(\epsilon)}\left[Y^{\epsilon}(i+1)-Y^{\epsilon}( i)\right]\right)\,.
			\end{split}
		\end{equation*}
		By using the invariance property of the measure $\pi^{\text{inv}}_{\alpha}$ we have that 
		\begin{equation*}
			\begin{split}
				N^{\epsilon}=
				\lfloor\epsilon^{-1}\rfloor\sum_{z\in \mathbb{Z}}\pi^{\text{inv}}_{\alpha}(z)\mathbb{COV}^{(\epsilon)}\left(\mathrm{E}_{z}^{(\epsilon)}\left[X^{\epsilon}(1)-z\right],\mathrm{E}_{0}^{(\epsilon)}\left[Y^{\epsilon}(1)\right]\right)\,.
			\end{split}
		\end{equation*}
		Therefore, it is enough to compute the covariance for one step of the chain, which we re-write as $\mathbb{COV}^{(\epsilon)}\left(\sum_{x\in \mathbb{Z}}(x-z)K_{1}^{\epsilon}(z,x),\sum_{y\in \mathbb{Z}}yK_{1}^{\epsilon}(0,y)\right)$. 

	If $|z|>1$, this covariance is trivially vanishing, then we consider the remaining cases. 
In what follows, for the sake of notation, we denote by $\bm{\mathcal{I}}$ a random vector containing the $\text{Bernoulli}(\epsilon)$ variables, associated with the environment and relevant for the considered case. 


\item{\underline{Case $z=0$.}} Here, we have that 
\begin{equation*}
	\mathbb{COV}^{(\epsilon)}\left(\sum_{x\in \mathbb{Z}}x K_{1}^{\epsilon}(0,x),\sum_{y\in \mathbb{Z}}yK_{1}^{\epsilon}(0,y)\right)=\mathbb{VAR}^{(\epsilon)}\left(\sum_{x\in \mathbb{Z}}xK_{1}^{\epsilon}(0,x)\right)\,.
\end{equation*}
The relevant $\text{Bernoulli}(\epsilon)$ variables are listed in  $\bm{\mathcal{I}}=(\mathcal{I}_{-1,0},\mathcal{I}_{0,1})\in \{0,1\}^{2}$. 
By the law of total variance we have that
\begin{equation*}
	\begin{split}
		&\mathbb{VAR}^{(\epsilon)}\left(\sum_{x\in \mathbb{Z}}xK_{1}^{\epsilon}(0,x)\right)
		\\=&\mathbb{E}^{(\epsilon)}\left[\mathbb{VAR}^{(\epsilon)}\left(\sum_{x\in \mathbb{Z}}xK_{1}^{\epsilon}(0,x)\bigg|\bm{\mathcal{I}}\right)\right]+\mathbb{VAR}^{(\epsilon)}\left(\mathbb{E}^{(\epsilon)}\left[\sum_{x\in \mathbb{Z}}xK_{1}^{\epsilon}(0,x)\bigg|\bm{\mathcal{I}}\right]\right)\,.
	\end{split}
\end{equation*}

For the first addend we obtain that 
\begin{equation*}
	\begin{split}
		\mathbb{E}^{(\epsilon)}\left[\mathbb{VAR}^{(\epsilon)}\left(\sum_{x\in \mathbb{Z}}xK_{1}^{\epsilon}(0,x)\bigg| \bm{\mathcal{I}}\right)\right]&= 
		2\epsilon \mathbb{VAR}^{(\epsilon)}(B_{0})+o(\epsilon)
		=
		\frac{\epsilon}{2(2\alpha+1)}+o(\epsilon)\,.
	\end{split}
\end{equation*}
{\color{black} The second addend gives
\begin{equation}
	\begin{split}
		\mathbb{VAR}^{(\epsilon)}\left(\mathbb{E}^{(\epsilon)}\left[\sum_{x}xK_{1}^{\epsilon}(0,x)\bigg|\mathbb{\mathcal{I}}\right]\right)=& \mathbb{VAR}^{(\epsilon)}\left(\mathbb{E}^{(\epsilon)}\left[B_{0,1}\mathcal{I}_{0,1}-(1-B_{-1,0})\mathcal{I}_{-1,0}\bigg|\mathbf{\mathcal{I}}\right]\right)
		\\=&\frac{1}{4}\mathbb{VAR}^{(\epsilon)}(\mathcal{I}_{0,1}-\mathcal{I}_{-1,0})
		\\=&\frac{1}{4}2\epsilon(1-\epsilon)=\frac{\epsilon}{2}+o(\epsilon)\,.
	\end{split}
\end{equation}
It follows that, summing the 2 contributions, we have
\begin{equation}
	\begin{split}
		\mathbb{VAR}^{(\epsilon)}\left(\sum_{x}xK_{1}^{\epsilon}(0,x)\right)= \frac{\epsilon}{2(2\alpha+1)}+\frac{\epsilon}{2}+o(\epsilon)= \frac{\alpha+1}{2\alpha+1}\epsilon+o(\epsilon)\,.
	\end{split}
\end{equation}}
\item{\underline{Case $z=1$.}} The relevant $\text{Bernoulli}(\epsilon)$ variables are listed in $\bm{\mathcal{I}}:=(\mathcal{I}_{-1,0},\mathcal{I}_{0,1},\mathcal{I}_{1,2})$. By applying the law of total covariances, we have that 
\begin{multline*}
		\mathbb{COV}^{(\epsilon)}\left(\sum_{x\in \mathbb{Z}}(x-1)K_{1}^{\epsilon}(1,x),\sum_{y\in \mathbb{Z}}yK_{1}^{\epsilon}(0,y)\right)
		=\\
		\mathbb{E}^{(\epsilon)}\left[\mathbb{COV}^{(\epsilon)}\left(\sum_{x\in \mathbb{Z}}(x-1)K_{1}^{\epsilon}(1,x),\sum_{y\in \mathbb{Z}}yK_{1}^{\epsilon}(0,y)\bigg|\bm{\mathcal{I}}\right)\right]
		\\+\mathbb{COV}^{(\epsilon)}\left(\mathbb{E}^{(\epsilon)}\left[\sum_{x\in \mathbb{Z}}(x-1)K_{1}^{\epsilon}(1,x)\bigg|\bm{\mathcal{I}}\right],\mathbb{E}^{(\epsilon)}\left[\sum_{y\in \mathbb{Z}}yK_{1}^{\epsilon}(0,y)\bigg|\bm{\mathcal{I}}\right]\right)\,.
\end{multline*}
The first addend is of order $\epsilon$ only when $\bm{\mathcal{I}}=(0,1,0)$, otherwise it is equal to zero or of order $o(\epsilon)$. Then, we write
\begin{equation}
	\begin{split}
		\mathbb{E}^{(\epsilon)}\left[\mathbb{COV}^{(\epsilon)}\left(\sum_{x\in \mathbb{Z}}(x-1)K_{1}^{\epsilon}(1,x),\sum_{y\in \mathbb{Z}}yK_{1}^{\epsilon}(0,y)\,\bigg|\,\bm{\mathcal{I}}\right)\right]=&
		\mathbb{VAR}^{(\epsilon)}(B_{0})\epsilon+o(\epsilon)
		\\=&
		\frac{\epsilon}{4(2\alpha+1)}+o(\epsilon)\,.
	\end{split}
\end{equation}
{\color{black}The second addend gives 
	\begin{equation}
		\begin{split}
			&\mathbb{COV}^{(\epsilon)}\left(\mathbb{E}^{(\epsilon)}\left[\sum_{x}(x-1)K_{1}^{\epsilon}(1,x)\bigg|\mathbf{\mathcal{I}}\right],\mathbb{E}^{(\epsilon)}\left[\sum_{y}yK_{1}^{\epsilon}(0,y)\bigg|\mathbf{\mathcal{I}}\right]\right) 
			\\=& \frac{1}{4}\mathbb{COV}^{(\epsilon)}\left(\mathcal{I}_{1,2}-\mathcal{I}_{0,1},\mathcal{I}_{0,1}-\mathcal{I}_{-1,0}\right)
			\\=&
			\frac{1}{4}\mathbb{E}^{(\epsilon)}\left[(\mathcal{I}_{1,2}-\mathcal{I}_{0,1})(\mathcal{I}_{0,1}-\mathcal{I}_{-1,0})\right]
			\\=& 
			\frac{1}{4}\left(\mathbb{E}^{\epsilon}[\mathcal{I}_{1,2}\mathcal{I}_{0,1}]+\mathbb{E}^{\epsilon}[\mathcal{I}_{0,1}\mathcal{I}_{-1,0}]-\mathbb{E}^{\epsilon}[\mathcal{I}_{1,2}\mathcal{I}_{-1,0}]-\mathbb{E}^{\epsilon}[\mathcal{I}_{0,1}^{2}]\right)
			\\=&
			\frac{1}{4}(\epsilon^{2}-\epsilon)\,.
		\end{split}
	\end{equation}
	It follows that 
	\begin{equation}\label{cov-z1}
		\mathbb{COV}^{(\epsilon)}\left(\sum_{x}(x-1)K_{1}^{\epsilon}(1,x),\sum_{y}yK_{1}^{\epsilon}(0,y)\right)= \frac{\eps
		}{4(2\alpha+1)}-\frac{1}{4}\epsilon+o(\epsilon)= -\frac{\alpha \epsilon}{2(2\alpha+1)}+o(\epsilon)
	\end{equation}}
\item{\underline{Case $z=-1$.}} By symmetry, one obtains again \eqref{cov-z1}. 

Using equation \eqref{N-epsilon} for $N^{\epsilon}$ and the measure $\pi^{\text{inv}}_{\alpha}$, written in \eqref{normalized-piInv}, we have that 
\begin{equation}\label{numerator-approx-final}
	N^{\epsilon}=\lfloor\epsilon^{-1}\rfloor
{\color{black}	\left(\epsilon \frac{(\alpha+1)}{\alpha}\frac{\alpha+1}{2\alpha+1}-\epsilon\frac{\alpha}{(2\alpha+1)}\right)+o(1)}
	=\frac{1}{\alpha}+o(1)\,.
\end{equation}

\bigskip 
Let us now turn to the denominator. 
		By using the telescopic identity \eqref{telescopic}, the invariance of the measure $\pi^{\text{inv}}_{\alpha}$ and by an argument similar to that done for the numerator, we obtain 
	\begin{equation}\label{denominator-approximation}
		D^{\epsilon}=\lfloor\epsilon^{-1}\rfloor
		\sum_{z\in \mathbb{Z}}\pi^{\text{inv}}_{\alpha}(z)\mathbb{E}^{(\epsilon)}\left[\mathrm{E}_{z}^{\text{dif},(\epsilon)}\left[|Z^{\epsilon}(1)|-|Z^{\epsilon}(0)|\right]\right]\,.
	\end{equation}
	Here, $\mathrm{E}_{z}^{\text{dif},(\epsilon)}$ denotes the quenched expectation with respect to the law of the Markov chain $(Z(n))_{n\in \mathbb{Z}_{\geq 0}}$ introduced in Definition \ref{remark-Zn}. Therefore, to find the exact expression for the denominator it is enough to compute the annealed expectation
	\begin{equation*}
		\begin{split}
			\mathbb{E}^{(\epsilon)}\left[\mathrm{E}_{z}^{\text{dif},(\epsilon)}\left[|Z^{\epsilon}(1)|\right]\right]-|z|=&\mathbb{E}^{(\epsilon)}\left[\mathbb{E}^{(\epsilon)}\left[\sum_{y\in \mathbb{Z}}K_{1}^{\epsilon}(0,y)\sum_{a\in \mathbb{Z}}K_{1}^{\epsilon}(z,a+y)|a|
			\bigg| \bm{\mathcal{I}}\right]\right]-|z|\,.
		\end{split}
	\end{equation*}
	For $|z|>1$, the denominator is vanishing. We then consider the remaining cases.
\item{\underline{Case $z=0$.}}  Here, the relevant $\text{Bernoulli}(\epsilon)$ variables are listed in $\bm{\mathcal{I}}=(\mathcal{I}_{-1,0},\mathcal{I}_{0,1})$. Then, we write
	\begin{equation*}
		\begin{split}
			\mathbb{E}^{(\epsilon)}\left[\mathbb{E}^{(\epsilon)}\left[\sum_{y\in \mathbb{Z}}K_{1}^{\epsilon}(0,y)\sum_{a\in \mathbb{Z}}K_{1}^{\epsilon}(0,a+y)|a|\bigg| \mathcal{I}\right]\right]=&
			\left(1-4\mathbb{VAR}^{(\epsilon)}(B)\right)\epsilon+o(\epsilon)
			\\=&
			\frac{2\alpha}{2\alpha+1}\epsilon+o(\epsilon)\,.
		\end{split}
	\end{equation*}
	
	\item{\underline{Case $z=1$.}} Here, the relevant $\text{Bernoulli}(\epsilon)$ variables are listed in  $\bm{\mathcal{I}}=(\mathcal{I}_{-1,0},\mathcal{I}_{0,1},\mathcal{I}_{1,2})$. 	 In this case, we have contribution of order at most $\epsilon$ only when $\mathcal{I}_{-1,0}=\mathcal{I}_{0,1}=\mathcal{I}_{1,2}=0$ (that is an event with probability $1-3\epsilon+o(\epsilon)$) or when  $\mathcal{I}_{-1,0}+\mathcal{I}_{0,1}+\mathcal{I}_{1,2}=1$ (that are events with probability $\epsilon+o(\epsilon)$). Then, we obtain
	\begin{equation}\label{z1-denominator}
		\begin{split}
			&\mathbb{E}^{(\epsilon)}\left[\mathbb{E}^{(\epsilon)}\left[\sum_{y\in \mathbb{Z}}K_{1}^{\epsilon}(0,y)\sum_{a\in \mathbb{Z}}K_{1}^{\epsilon}(1,a+y)|a|\bigg| \bm{\mathcal{I}}\right]\right]-1
\\=&
			  \left(\left(1+\mathbb{E}^{(\epsilon)}[B_{-1}]\right)+2\left(\frac{1}{4}-\mathbb{VAR}^{(\epsilon)}(B_{0})\right)
			  +\mathbb{E}^{(\epsilon)}[2(1-B_{1})+B_{1}]\right)\epsilon-3\epsilon+o(\epsilon)
\\=&
		\frac{\alpha}{2\alpha+1}\epsilon+o(\epsilon)
		\end{split}
	\end{equation}
\item{\underline{Case $z=-1$.}}  By symmetry, we have again \eqref{z1-denominator}.

	By computing the expectation with respect to the invariant measure $\pi^{\text{inv}}_{\alpha}$ of equation \eqref{normalized-piInv}, we obtain 
	\begin{equation}\label{denominator-approx-final}
		\begin{split}
			D^{\epsilon}=&\lfloor\epsilon^{-1}\rfloor\left(\frac{2\alpha}{2\alpha+1}\epsilon+\frac{2\alpha}{2\alpha+1}\frac{\alpha+1}{\alpha}\epsilon\right)+o(1)=2+o(1)\,.
		\end{split}
	\end{equation}
	Finally, by using \eqref{numerator-approx-final} and \eqref{denominator-approx-final}, we have that
	\begin{equation*}
		\lim_{\epsilon\to 0}\gamma_{\epsilon}^{2}=  \lim_{\epsilon\to 0}\frac{N^{\epsilon}}{D^{\epsilon}}=\frac{1}{{\color{black}2}\alpha}\,.
	\end{equation*}
\end{proof}

\end{document}
